\documentclass[11pt,a4paper]{article}
\usepackage[T1]{fontenc}
\usepackage{lmodern}
\usepackage[margin=21mm]{geometry}
\usepackage{amsmath,amssymb,booktabs,tabularx,array}
\usepackage[hidelinks]{hyperref}
\usepackage{microtype}
\hypersetup{pdftitle={Two notions of polynomial time on the word-RAM},
  pdfsubject={Lax759944, Lax560851, input magnitude, and word capacity}}
\setlength{\parindent}{0pt}
\setlength{\parskip}{6pt}
\newcolumntype{Y}{>{\raggedright\arraybackslash}X}
\newcommand{\code}[1]{\texttt{#1}}
\newcommand{\bits}{\operatorname{bitSize}}
\title{Two notions of polynomial time on the word-RAM}
\author{}
\date{Lax759944 and the current Lax560851 example --- 16 September 2026}

\begin{document}
\begin{center}
  {\Large\bfseries Two notions of polynomial time on the word-RAM}\par
  {\small Lax759944 and the current Lax560851 example\enspace\textbullet\enspace 16 September 2026}
\end{center}

The definitions differ in their input-size measure and required word-width
threshold. This note compares \code{RamPolytime} in Lax759944 with Lax560851's
\code{PolynomialTime} example. Both use Lax808846: read-only input, indexed
and sequential access, zero initial working memory, and append-only output.
Fetched terminal instructions count toward time.

For a list $x$ of natural numbers, use the following quantities:
\[
\begin{aligned}
  n &= \text{number of entries in }x,\\[3pt]
  m &= \text{largest entry in }x\text{ (zero for the empty list)},\\[3pt]
  B &= \bits(x)=\text{length of Lax759944's binary encoding},\\[3pt]
  M &= \operatorname{inputMagnitude}(x)=2n+m+2.
\end{aligned}
\]
The binary encoding has one separator per entry, followed by that entry's
binary digits. Thus $B\ge n$. For example, $[2^{1000}]$ has $n=1$ but $B=1002$.

\section*{1. Definitions and units}

\begingroup
\small
\renewcommand{\arraystretch}{1.16}
\begin{tabularx}{\textwidth}{@{}p{0.20\textwidth}YY@{}}
\toprule
 & \textbf{Lax759944: RamPolytime} & \textbf{Lax560851: current PolynomialTime example} \\
\midrule
Time measure & Binary input size $B$. & Number of list entries $n$. \\
Time allowance & $P(B)$ for a polynomial $P$. & $A(n+1)^d$ for constants $A,d\in\mathbb N$. \\
Word requirement & $w\ge Q(B)$, where $Q$ is a polynomial. The bound is in \emph{bits}. & $CM\le 2^w$, plus payload and address fit, for a constant $C$. The bound is in \emph{numerical capacity}. \\
Input magnitude $M$ & Not used as the size measure; size is the binary encoding length $B$. & $M=S+m+1$: structural node count $S$, plus largest natural payload $m$, plus one. Here $S=2n+1$, so $M=2n+m+2$. This is neither bit length nor memory usage. \\
Physical input & The native list with its length prepended: $n::x$. & The certified structural arena: a root word and three words per structural node. \\
Output compared here & The singleton list $[f(x)]$. & The same singleton list $[f(x)]$, because $f$ returns a natural number. \\
Uniformity & One program, working at every sufficiently large width. & The same quantifier discipline. \\
\bottomrule
\end{tabularx}
\endgroup

Neither definition imposes an upper limit on the actual word width $w$.
Each guarantees a run at \emph{every} width above its sufficient threshold.
The question is how small that threshold must be, not how large the machine
is allowed to be.

Lax759944 requires the native input and exact output to fit at width $Q(B)$,
and then requires a run of at most $P(B)$ instructions for every $w\ge Q(B)$.
For example, the sufficient width could be $B^2$ bits, with numerical
capacity $2^{B^2}$.

For a list of naturals, Lax560851's three width premises are exactly
\[
  m<2^w,\qquad 6n+3\le 2^w,\qquad C(2n+m+2)\le 2^w.
\]
The middle condition addresses the arena's $3S=6n+3$ memory words.
For fixed $C$, all three conditions hold starting at a width of order
$\log(n+m+2)$. For bounded payloads, this is $O(\log n)$ bits.

\newpage
\section*{2. Why input magnitude contains $m$, not $\log m$}

A $w$-bit word represents integers from $0$ through $2^w-1$. Storing the
value $m$ therefore requires $m<2^w$, or about $\log_2(m+1)$ bits. The
logarithm appears when solving for the \emph{width}; it does not belong on
the capacity side of the inequality.

Accordingly, $C(2n+m+2)\le 2^w$ already gives logarithmic dependence of
the required width on $m$. Replacing $m$ by $\log m$ on the left would mix
the bit length of a value with the numerical range needed to represent it.
The separate payload-fit premise would still be necessary. Input magnitude
is a convenient combined scale for address range and payload values, not a
claim that either the storage or the number of bits is proportional to $m$.

\section*{3. A function separating the definitions}

Consider
\[
  f(x)=2^{\operatorname{length}(x)}.
\]
For Lax759944, use the output adapter $x\mapsto[f(x)]$ to match the natural-number
output of the Lax560851 example. The value $2^n$ needs $n+1$ bits, even though its
numerical value is exponential in $n$.

\textbf{Why Lax759944 accommodates it.}
The physical input starts with $n$. A fixed RAM program reads this prefix,
sets a register to $1$, shifts it left by $n$ bits, writes the result, and
halts. It need not inspect the remaining input entries. Left shift is an
instruction in the underlying RAM model, so the instruction count is constant.
Since $n\le B$, the polynomial sufficient width $Q(B)=B+2$ accommodates
the input entries, the length prefix, and the output. The program works at
every larger width as well.

\textbf{Why the current Lax560851 example excludes it.}
Let $x$ consist of $n$ zeros. Then $m=0$, $B=n$, and $M=2n+2$.
For any proposed constant $C$, choose the least $w$ for which
\[
  2^w\ge \max\{6n+3,\ C(2n+2)\}.
\]
This width satisfies all three Lax560851 premises and is $O(\log n)$.
Nevertheless, its required one-word output $2^n$ fits only if $w\ge n+1$.
For sufficiently large $n$, these conditions are incompatible. Increasing
the time coefficient or degree cannot fix an output that does not fit in a
word. Returning several words would change the specified output encoding.

For illustration, take $C=10$ and $n=1000$: the premises hold at $w=15$,
because $2^{15}=32768\ge20020$. But the output requires $1001$ bits.
For any other fixed $C$, a sufficiently large $n$ produces the same obstruction.

\section*{4. Consequence and possible alignment}

The definitions are not equivalent, even on zero lists where $B=n$. The
general Lax560851 predicate is flexible enough to use time bound $P(B)$ and
capacity bound $2^{Q(B)}$; its capacity premise then becomes $Q(B)\le w$.
The arena-based \code{BitPolynomialTime} definition makes this alignment.
Its equivalence with Lax759944 is proved for all Lax808846 instructions.
The verified conversions preserve indexed input, sequential reads, length,
EOF, and exact output, and may add linear startup time in either direction.
They choose separate polynomial witnesses; no equality of degrees is required.

\small
\textit{Formalization, 16 September 2026.} The output-fit invariant,
separation theorem, seven-instruction exponential example, and general
equivalence are checked in Lean 4.33. The equivalence uses the verified
compiler for the complete Lax808846 instruction set, with input initialization
and terminal instructions included in the time bounds.
\end{document}
